\begin{abstract}
End-to-end driving policies are evaluated on benchmarks whose images are, almost without
exception, clear and well-lit. We ask what their internal representations do when only the
\emph{appearance} of a scene changes, and find a failure mode that is consistent across four
architectures: a synthetic night rendering moves the representation along a direction that is
systematically \emph{opposed} to the policy's own hazard-response direction, and the planned
speed rises accordingly. In representation space the hazard direction and the braking
direction form one family ($34$--$46^\circ$ apart, against a $7^\circ$ reproducibility floor
and a $90^\circ$ random baseline), while the appearance direction sits past the random line on
the far side ($93$--$115^\circ$). Behaviourally, the same perturbation raises planned speed by
up to $2.05$\,m/s --- an order of magnitude more than deleting the lead vehicle from the image.
Two controls sharpen the mechanism. First, the effect is carried by sensor noise, not by low
luminance: darkening alone changes planned speed by $-0.01$\,m/s, and adding noise of the
magnitude a real night camera produces flips it to $+0.05$. Second, perturbing \emph{only the
sky} --- $17.9\%$ of the frame, with the driving corridor untouched to the pixel --- still
recovers $19\%$ of the effect of removing the hazard itself, so the dependence is global rather
than mediated by any local cue. We also report, from five pre-registered tests, that the
projection onto these directions is \emph{not} a reliable proxy for per-policy or per-event
behaviour; four of the five are rejected. Direction readouts are widely used as behavioural
proxies, and we give multi-angle evidence for where that use breaks down.
\end{abstract}
